\slajdid{09_2-s048}
\begin{frame}[label=09_2-s048]
\gusttekst\setlength{\abovedisplayskip}{3pt}\setlength{\belowdisplayskip}{3pt}\setlength{\abovedisplayshortskip}{2pt}\setlength{\belowdisplayshortskip}{2pt}Zbog mogućnosti skaliranja i prelazaka sa jedne na druge tehnologije uveden je normirana dimenzija, parametar $\lambda$ (nema veze sa parametrima $\lambda_n$ i $\lambda_n$) koji u stvari pokazuje mogućnosti tehnologije, a svi ostale dimenzije koji se daju su celobrojni umnošci tog parametra. Znači dimenzije tranzistora i odnosi dimenzija se daju samo kao celobrojne vrednosti. Parametar $\lambda$ je definisan tako da $2\lambda$ u stvari predstavlja minimalne tehnološke dimenzije. Neke dimenzije je moguće ostvariti i sa dimenzijama koje su „polovina minimalne dimenzije“. Prilikom označavanja tranzistora najčešće se podrazumeva da je dužina kanala svih tranzistora ista (i često minimalna) a onda se označavaju samo širine kanala.\par\medskip
\begin{columns}[c,onlytextwidth]\column{.21\textwidth}\centering\input{slike/tikz/s048_cmos_invertor.tex}\column{.77\textwidth}
U ovom smislu se i definiše \alert{\textbf{CMOS invertor minimalne geometrije – jedinični CMOS invertor}}. Dužine kanala P i N tranzistora su iste, a širine zadovoljavaju neke od prethodnih izvedenih odnosa, ili kompromis tih zahteva. Najčešće je odnos 2:1. Zauzima „najmanje prostora“ i kapacitivnosti su minimalne.\par\bigskip
U ovom smislu treba biti samo oprezan, kada se daju odnosi „većih tranzistora“. Na primer neka je računski na osnovu nekih prethodnih uslova dobijeno da odnos treba da bude 2.376. Za invertor minimalne geometrije taj odnos treba da bude približan ceo broj odnosno 2. Međutim ako je iz nekog razloga kanal N tranzistora širine 10, onda P tranzistor neće biti širine 20 nego 24, itd…
\end{columns}
\end{frame}
